\section{Finite shape bounds}\label{sec:convex}

This section proves the finite bounds behind
Theorems~\ref{thm:main-limit} and~\ref{thm:two-modes}. For a fair
start, Theorem~\ref{thm:five-blocks} shows that the second
difference of the count law changes sign at most four times. For
every initial bias, Proposition~\ref{prop:all-bias-modes} shows
that the count law has at most two modal plateaus. The limit
transfer is Lemma~\ref{lem:derivative-limits} in
Appendix~\ref{app:modes}.

We first treat second differences under a fair start. Let
$0<\varepsilon<\frac12$, $a=1-2\varepsilon$, and $r_t(k)=\varepsilon+ak/t$. For $t\ge1$ and every integer $k$ we write
$P_t(k)=\Prob(A_t=k)$, which is $0$ unless $0\le k\le t$, and
\[
 D_t(k)=P_t(k+1)-2P_t(k)+P_t(k-1).
\]
We also put $s_t=D_t(1)=P_t(0)-2P_t(1)+P_t(2)$ for $t\ge2$.

\begin{lemma}[recurrence for the second difference]\label{lem:D-rec}
For every $t\ge1$ and every integer $k$,
\begin{align}
 P_{t+1}(k)&=\bigl(1-r_t(k)\bigr)P_t(k)+r_t(k-1)\,P_t(k-1),\label{eq:P-rec}\\
 D_{t+1}(k)&=\bigl(1-r_t(k+1)\bigr)D_t(k)+r_t(k-2)\,D_t(k-1).\label{eq:D-rec}
\end{align}
Also $P_t(k)=P_t(t-k)$ and $D_t(k)=D_t(t-k)$.
\end{lemma}

\begin{proof}
Given $A_t=k$, the copied step is $+1$ with probability $k/t$, so step $t+1$
is $+1$ with probability $(1-\varepsilon)\frac kt+\varepsilon(1-\frac
kt)=r_t(k)$. This gives \eqref{eq:P-rec} for $0\le k\le t+1$. At $k=0$ and
$k=t+1$, the extra term $r_t(-1)P_t(-1)$ or
$(1-r_t(t+1))P_t(t+1)$ is $0$ because $P_t(-1)=P_t(t+1)=0$.
Outside $0\le k\le t+1$ both sides are $0$.
Under a fair start, $(X_i)$ and $(-X_i)$ have the same law, so
$P_t(k)=P_t(t-k)$. Taking second differences gives the same for $D_t$.

For \eqref{eq:D-rec}, take second differences in \eqref{eq:P-rec}.
The affine identity $r_t(k+1)-r_t(k)=a/t$ gives
\[
 2r_t(k+1)+r_t(k-2)=3r_t(k),\qquad
 r_t(k+1)+2r_t(k-2)=3r_t(k-1),
\]
so the coefficients of $P_t(k)$ and $P_t(k-1)$ agree with the
right side of~\eqref{eq:D-rec}. The outer two coefficients agree directly.
\end{proof}

Trinas and Valle~\cite[(2.2)]{tv2026} use \eqref{eq:P-rec}, the forward equation, for a superdiffusive local
limit theorem. Their higher-difference
recurrence~\cite[(4.10)]{tv2026} is the basis of a computer-certified
study of log-concavity at fixed memory. We use \eqref{eq:D-rec} for a different
purpose, to count sign changes. For $2\le k\le t-1$ its coefficients
are positive. So new changes can only come in at the ends.
Lemma~\ref{lem:interleave} shows that positive sums of neighbors never add a change
(see~\cite{karlin1968,karlin1964} for the general theory).

Let $z$ be a finite real list. Delete its zeros, replace each entry
left by its sign, and merge equal neighbors. What remains is the
sign pattern $\operatorname{sp}(z)$, an alternating list of signs. For
two such lists we write $\sigma\preceq\tau$ if $\sigma$ is a subsequence of $\tau$. We put
$\Pi_5=(+,-,+,-,+)$ and $\Pi_3=(+,-,+)$.

\begin{lemma}[positive sums of neighbors]\label{lem:interleave}
Let $\zeta=(\zeta_1,\dots,\zeta_n)$ be real, and let
$\eta_i=A_i\zeta_i+B_i\zeta_{i+1}$ for $1\le i\le n-1$ with $A_i\ge0$ and
$B_i\ge0$. Let $w=(\zeta_1,\eta_1,\zeta_2,\eta_2,\dots,\eta_{n-1},\zeta_n)$.
Then $\operatorname{sp}(w)=\operatorname{sp}(\zeta)$ and
$\operatorname{sp}(\eta)\preceq\operatorname{sp}(\zeta)$.
\end{lemma}

\begin{proof}
Insert each $\eta_i$ between its two neighbors in $\zeta$. If $\eta_i=0$, the
sign pattern does not change. If $\eta_i>0$, then $\zeta_i>0$ or
$\zeta_{i+1}>0$, so $\eta_i$ is placed next to an entry of its own sign, and
the pattern does not change. The same holds for $\eta_i<0$, proving the first claim.
For the second, obtain $\eta$ from $w$ by deleting the entries
$\zeta_i$. Deletion can only remove signs from the pattern.
\end{proof}

For the walk we apply Lemma~\ref{lem:interleave} to the augmented
rows of Figure~\ref{fig:augmented-rows}, drawn there for first
differences. Proposition~\ref{prop:all-bias-modes} treats the whole
row and Proposition~\ref{prop:dominant-half} the half $2k+1\ge t$.

\begin{figure}[ht]
\centering
\small
\[
\begin{array}{c@{\quad}c@{\quad}c}
\underbrace{\ell_N,\ldots,\ell_{t+1}}_{\text{left boundary history}}&
\underbrace{d_t(0),\ldots,d_t(t-1)}_{\text{active row}}&
\underbrace{u_{t+1},\ldots,u_N}_{\text{right boundary history}}\\[9pt]
&\downarrow\ \text{copy or sum adjacent entries}&\\[9pt]
\ell_N,\ldots,\ell_{t+2}&
\ell_{t+1},d_{t+1}(1),\ldots,d_{t+1}(t-1),u_{t+1}&
u_{t+2},\ldots,u_N
\end{array}
\]
\caption{One step of the augmented rows for a final length $N$, with
$d_t(k)=P_t(k+1)-P_t(k)$, $\ell_t=d_t(0)$, and $u_t=d_t(t-1)$. The
active row gains an entry and the histories on each side lose one.
Each new end value was already in a history, so the bound for the
previous row already counts its sign.}
\label{fig:augmented-rows}
\end{figure}

\begin{lemma}[sign changes next to an end]\label{lem:level-one}
For every $t_1\ge2$ we have $\operatorname{sp}(s_2,s_3,\dots,s_{t_1})\preceq\Pi_3$.
\end{lemma}

We prove this in Appendix~\ref{app:signs-level}.

\begin{theorem}[at most five blocks]\label{thm:five-blocks}
Let $N\ge3$ and $0<\varepsilon<\frac12$. Then
\[
 \operatorname{sp}\bigl(D_N(1),D_N(2),\dots,D_N(N-1)\bigr)\preceq(+,-,+,-,+).
\]
So the second difference of $P_N$ changes sign at most four times on
$\{1,\dots,N-1\}$.
\end{theorem}

\begin{proof}
For $3\le t\le N$, we retain the endpoint history in the list $\mathcal Z_t$
\[
 \mathcal Z_t=\bigl(s_N,s_{N-1},\dots,s_{t+1},\ s_t,D_t(2),D_t(3),\dots,D_t(t-2),s_t,\
 s_{t+1},\dots,s_{N-1},s_N\bigr).
\]
It has $2N-t-1$ entries. Its middle part is
$(D_t(1),D_t(2),\dots,D_t(t-1))$, because $D_t(t-1)=D_t(1)=s_t$ by symmetry.
The list $\mathcal Z_3$ is $(s_N,\dots,s_4,s_3,s_3,s_4,\dots,s_N)$, and $\mathcal Z_N$ is the
list in the statement.

To apply Lemma~\ref{lem:interleave}, we check how
$\mathcal Z_{t+1}$ is obtained from $\mathcal Z_t$ for $3\le t\le N-1$. The first $N-t$ entries of
$\mathcal Z_{t+1}$ are $s_N,\dots,s_{t+1}$, and they are copies of the first $N-t$
entries of $\mathcal Z_t$. So $A_i=1$ and $B_i=0$ there. The next $t-2$ entries are
$D_{t+1}(j)$ for $2\le j\le t-1$. By \eqref{eq:D-rec},
\[
 D_{t+1}(j)=r_t(j-2)\,D_t(j-1)+\bigl(1-r_t(j+1)\bigr)D_t(j),
\]
and $D_t(j-1)$ and $D_t(j)$ are neighbors in the middle part of $\mathcal Z_t$, in the
required order. The coefficients satisfy $r_t(j-2)\ge\varepsilon$ and
$1-r_t(j+1)\ge1-\varepsilon-a=\varepsilon$, since $2\le j\le t-1$. The last
$N-t$ entries of $\mathcal Z_{t+1}$ are $s_{t+1},\dots,s_N$, and each is a copy of the
entry of $\mathcal Z_t$ one place to its right. So $A_i=0$ and $B_i=1$ there.

By Lemma~\ref{lem:interleave},
$\operatorname{sp}(\mathcal Z_N)\preceq\operatorname{sp}(\mathcal Z_{N-1})\preceq\dots\preceq\operatorname{sp}(\mathcal Z_3)$.
By Lemma~\ref{lem:level-one} the sign pattern $p$ of $(s_3,\dots,s_N)$ is a
subsequence of $\Pi_3$, so it is one of the lists $()$, $(+)$, $(-)$, $(+,-)$,
$(-,+)$, and $(+,-,+)$. The sign pattern of $\mathcal Z_3$ is $p$ reversed, followed by
$p$, with equal neighbors merged. In the six cases this is $()$, $(+)$, $(-)$,
$(-,+,-)$, $(+,-,+)$, and $(+,-,+,-,+)$, and each is a subsequence of $\Pi_5$.
\end{proof}

\subsection{First differences for a biased start}\label{sec:biased-finite}

\begin{proposition}[two modes for every initial bias]\label{prop:all-bias-modes}
Let $0<\varepsilon<1$, let $q\in[0,1]$ be the probability that the first
step is positive, and let $P_n(k)$ be the law of the number of positive
steps after $n\ge1$ steps. Extend $P_n$ by zero outside $\{0,\ldots,n\}$.
The first differences $d_n(k)=P_n(k+1)-P_n(k)$ satisfy
\[
 \operatorname{sp}\bigl(d_n(-1),d_n(0),\ldots,d_n(n)\bigr)
 \preceq (+,-,+,-).
\]
In particular, $P_n$ has at most two modal plateaus, including endpoint
modes. Moreover, there are no integers $i_0<\cdots<i_4$ for which
\[
 P_n(i_0)>P_n(i_1)<P_n(i_2)>P_n(i_3)<P_n(i_4).
\]
\end{proposition}

\begin{proof}
Put $a=1-2\varepsilon$ and $r_t(k)=\varepsilon+ak/t$. The forward equation
is
\[
 P_{t+1}(k)=(1-r_t(k))P_t(k)+r_t(k-1)P_t(k-1),
 \qquad P_1=(1-q,q).
\]
Subtracting adjacent equations gives
\begin{equation}\label{eq:bias-first-diff}
 d_{t+1}(k)=r_t(k-1)d_t(k-1)+(1-r_t(k+1))d_t(k),
 \qquad 1\le k\le t-1.
\end{equation}
Both coefficients are at least $\min(\varepsilon,1-\varepsilon)>0$,
since $r_t(j)$ is a convex combination of $\varepsilon$ and
$1-\varepsilon$ for $0\le j\le t$. It remains to check the two ends.

Suppose first that $q<1$. Since
$P_t(0)=(1-q)(1-\varepsilon)^{t-1}>0$, define
\[
 \xi_t=\frac{P_t(1)}{P_t(0)},\qquad
 b=\frac{1-2\varepsilon}{1-\varepsilon}<1,\qquad
 A=\frac{\varepsilon}{1-\varepsilon}>0.
\]
The boundary equation gives
\[
 \xi_{t+1}=(1-b/t)\xi_t+A,\qquad \xi_1=\frac q{1-q}.
\]
Thus $\delta_t=\xi_{t+1}-\xi_t$ obeys
\begin{equation}\label{eq:bias-edge-increments}
 \delta_{t+1}=\frac{t-b}{t+1}\delta_t+\frac A{t+1}.
\end{equation}
Once an increment is nonnegative, every later increment is positive.
Thus $\xi_t$ first decreases and then increases (either part may be empty). For $\ell_t=d_t(0)=P_t(0)(\xi_t-1)$ this proves
\[
 \operatorname{sp}(\ell_1,\ldots,\ell_N)\preceq(+,-,+).
\]
If $q\le1/2$, then $\xi_1\le1$ and the sharper bound is
$\operatorname{sp}(\ell_1,\ldots,\ell_N)\preceq(-,+)$.
In particular, a positive $\ell_i$ cannot be followed by a negative
$\ell_j$. Zeros do not affect either conclusion.

Write $u_t=d_t(t-1)$. Reflection of the walk gives
$u_t(q)=-\ell_t(1-q)$ and
$d_t^{(1-q)}(t-1-k)=-d_t^{(q)}(k)$.
For $0<q\le1/2$, the right boundary history therefore has pattern
contained in $(-,+,-)$. For $q=0$ it is nonpositive throughout. The conditional count is at
most $t-1$, so $u_t=-P_t(t-1)\le0$.

We put the values arriving at the boundary in an augmented row, since
they are not controlled by the interior recurrence. Fix $N$ and, for
$1\le t\le N$, form
\[
 Z_t=(\ell_N,\ldots,\ell_{t+1},
       d_t(0),\ldots,d_t(t-1),
       u_{t+1},\ldots,u_N).
\]
The left history is read backwards while the right is read forwards,
so every endpoint that arrives was already in the previous augmented
row and the recurrence ends up controlling the entire row.
For $q\le1/2$, its first row is contained in $(+,-,+,-)$.
We must also allow zeros at $\ell_1=u_1=2q-1$. Suppose $Z_1$
contained four ordered signs $-,+,-,+$. If the second lay in the
reversed left history, the first two would give
a positive left edge followed in time by a negative one, which is
impossible. Otherwise the last three lie in the right history and have
signs $+,-,+$. For $q>0$, their reflected negatives give the forbidden pattern
$(-,+,-)$ in the left history at bias $1-q$. For $q=0$, they contradict
the nonpositive right history. This also covers $N=1$.

Every entry of $Z_{t+1}$ is a nonnegative combination of adjacent entries
of $Z_t$. The interior entries use \eqref{eq:bias-first-diff}. The two new
boundary entries and the remaining future edge values are copied,
with coefficients $(1,0)$ or $(0,1)$. To check the ordering, number the augmented entries from left to
right. Output entry $i$ uses input entries $i$ and $i+1$, with a zero
coefficient when copied. A nonzero output
has a parent of the same sign. Two neighbors in the subsequence with
opposite signs cannot share a parent, since one entry has only one
sign. So the parent indices strictly increase. Pick parents like this
for an alternating subsequence of the output. We get the same signs,
in the same order, in the input.
An update therefore cannot add an alternating pattern, so
$\operatorname{sp}(Z_N)\preceq(+,-,+,-)$. Reflection and sign
reversal preserve this pattern, proving the claim also for $q\ge1/2$,
including $q=1$.

Finally, extension by zero adds $d_N(-1)=P_N(0)\ge0$ on the left and $d_N(N)=-P_N(N)\le0$ on the right. Neither enlarges the pattern.
Every modal plateau contributes a positive-to-negative change, so there
are at most two. If the statement's last inequalities held, telescoping
on each of their four intervals would give an adjacent difference
of the same strict sign. These four differences would
have signs $-,+,-,+$, a contradiction.
\end{proof}
